Última prueba de la parte: dibujar el paper y la réplica según el porcentaje de máquina ocupada, para comparar formas.

Con diez pilas y pocas escrituras, se parece: Autores encabeza en los dos. Con un millón de pilas, la forma se parece, aunque cambia el orden de arriba. Con una pila y pocas escrituras, no se parece: en el paper Anderson sube y adelanta, y aquí se queda plano.

Entonces, ¿simulamos los cuatro Xeon? Solo en parte. Es como un simulador de vuelo de juguete: la física que enseña es real y medida, pero el castigo por distancia es mucho más suave. Sirve para ver qué algoritmos son sensibles a la distancia, no para predecir los números del paper.
